\subsection{两个无穷和的乘积}

命题 1. 若有限实数的族 \((x_{\lambda})_{\lambda \in L}\) 与 \((y_{\mu})_{\mu \in M}\) 在 \(\mathbf{R}\) 中可和，则族 \((x_{\lambda}y_{\mu})_{(\lambda ,\mu)\in L\times M}\) 也可和，且我们有

\[
\sum_ {(\lambda , \mu) \in \mathbf {L} \times \mathbf {M}} x _ {\lambda} y _ {\mu} = (\sum_ {\lambda \in \mathbf {L}} x _ {\lambda}) (\sum_ {\mu \in \mathbf {M}} y _ {\mu}).\tag{1}
\]

\(\mathbf{L} \times \mathbf{M}\) 的每个有限子集都含于形如 \(\mathbf{H} \times \mathbf{K}\) 的有限子集，其中 \(\mathbf{H}\) 是 \(\mathbf{L}\) 的有限子集，\(\mathbf{K}\) 是 \(\mathbf{M}\) 的有限子集。根据假设，存在数 \(a > 0\) ，使得 \(\sum_{\lambda \in \mathbf{H}} |x_{\lambda}| \leqslant a\) 与 \(\sum_{\mu \in \mathbf{K}} |y_{\mu}| \leqslant a\) 对 \(\mathbf{H}\) 与 \(\mathbf{K}\) （\(\mathbf{L}\) 与 \(\mathbf{M}\) 的所有有限子集）分别成立；因此

\[
\sum_ {(\lambda , \mu) \in \mathbf {H} \times \mathbf {K}} | x _ {\lambda} y _ {\mu} | = (\sum_ {\lambda \in \mathbf {H}} | x _ {\lambda} |) (\sum_ {\mu \in \mathbf {K}} | y _ {\mu} |) \leqslant a ^ {2},
\]

而由定理 1 与定理 3，这表明族 \((x_{\lambda}y_{\mu})\) 在 \(\mathbf{R}\) 中可和。由结合律我们可以写{[}第三章，§ 5，第 3 节，公式 (2){]}

\[
\sum_ {(\lambda , \mu) \in \mathbf {L} \times \mathbf {M}} x _ {\lambda} y _ {\mu} = \sum_ {\lambda \in \mathbf {L}} (\sum_ {\mu \in \mathbf {M}} x _ {\lambda} y _ {\mu}) = \sum_ {\lambda \in \mathbf {L}} x _ {\lambda} (\sum_ {\mu \in \mathbf {M}} y _ {\mu}) = (\sum_ {\lambda \in \mathbf {L}} x _ {\lambda}) (\sum_ {\mu \in \mathbf {M}} y _ {\mu});
\]

由此即得结论。

